\paragraph{Solution for \cref{box:LOalphas}}
\begin{align}
    \alpha_s(\mu_R^2) &= \frac{\alpha_s(\mu_{R,0}^2)}{1 + \alpha_s(\mu_{R,0}^2)\beta_0\ln(\mu_R^2/\mu_{R,0}^2)}\\
    &= \frac 1 {\beta_0} \frac{1}{1/(\alpha_s(\mu_{R,0}^2)\beta_0) + \ln(\mu_R^2/\mu_{R,0}^2)}\\
    &= \frac 1 {\beta_0} \left\{\ln\left[\mu_R^2/\mu_{R,0}^2\exp(1/(\alpha_s(\mu_{R,0}^2)\beta_0))\right]\right\}^{-1}\\
    &= \frac 1 {\beta_0 \ln(\mu_R^2/\LQCD^2)}
\end{align}
with
\begin{equation}
    \LQCD^2 = \mu_{R,0}^2\exp(-1/(\alpha_s(\mu_{R,0}^2)\beta_0))
\end{equation}

\paragraph{Solution for \cref{box:LOalphassol}}
\begin{align}
    \dv {\alpha_s^\text{LO}(t_R)}{t_R} &= \dv {t_R}  \frac{\alpha_s(0)}{1 + \alpha_s(0) \beta_0 t_R}\\
     &= \frac{\alpha_s(0)}{(1 + \alpha_s(0) \beta_0 t_R)^2} \cdot (-1) \cdot \alpha_s(0) \beta_0\\
     &= -\beta_0 \left(\alpha_s^\text{LO}(t_R)\right)^2
\end{align}

\paragraph{Solution for \cref{box:NLOalphas}}
Step 1: compute LHS.
\begin{align}
    \dv {\alpha_s^\text{NLO}(t_R)}{t_R} &= \dv {t_R} \left(\alpha_s^\text{LO}(t_R) - \frac{\beta_1}{\beta_0} \left(\alpha_s^\text{LO}(t_R)\right)^2 \ln(\alpha_s(0) / \alpha_s^\text{LO}(t_R)) \right)\\
     &= \dv {\alpha_s^\text{LO}(t_R)}{t_R} - \frac{\beta_1}{\beta_0} \left[ 2 \alpha_s^\text{LO}(t_R) \dv {\alpha_s^\text{LO}(t_R)}{t_R} \ln(\alpha_s(0) / \alpha_s^\text{LO}(t_R)) \right. \nonumber\\
     &\hspace{120pt} \left. - \left(\alpha_s^\text{LO}(t_R)\right)^2 \dv{\ln(\alpha_s^\text{LO}(t_R))}{t_R}  \right]\\
     &= -\beta_0 \left(\alpha_s^\text{LO}(t_R)\right)^2 - \beta_1\left[  -2\left(\alpha_s^\text{LO}(t_R)\right)^3\ln(\alpha_s(0) / \alpha_s^\text{LO}(t_R)) + \left(\alpha_s^\text{LO}(t_R)\right)^3 \right]\\
     &= -\beta_0 \left(\alpha_s^\text{LO}(t_R)\right)^2 - \beta_1\left(\alpha_s^\text{LO}(t_R)\right)^3\left[1  -2\ln(\alpha_s(0) / \alpha_s^\text{LO}(t_R))\right]
\end{align}

Step 2: compute RHS.
\begin{align}
    &-\beta_0 \left(\alpha_s^\text{NLO}(t_R)\right)^2 -\beta_1 \left(\alpha_s^\text{NLO}(t_R)\right)^3 \nonumber\\
    &=-\beta_0 \left[\left(\alpha_s^\text{LO}(t_R)\right)^2 - 2 \frac{\beta_1}{\beta_0} \left(\alpha_s^\text{LO}(t_R)\right)^3 \ln(\alpha_s(0) / \alpha_s^\text{LO}(t_R)) \right] \nonumber\\
    &\hspace{20pt} - \beta_1 \left(\alpha_s^\text{LO}(t_R)\right)^3 + \order{(\alpha_s^\text{LO}(t_R))^4}\\
    &= -\beta_0 \left(\alpha_s^\text{LO}(t_R)\right)^2 - \beta_1\left(\alpha_s^\text{LO}(t_R)\right)^3\left[1  -2\ln(\alpha_s(0) / \alpha_s^\text{LO}(t_R))\right]  + \order{(\alpha_s^\text{LO}(t_R))^4}
\end{align}
where we immediately neglect all term $\order{(\alpha_s^\text{LO}(t_R))^4}$.

\paragraph{Solution for \cref{box:NLOalphasLL}}
Step 1: expand $\alpha_s^\text{LO}(t_R)$.
\begin{equation}
    \alpha_s^\text{LO}(t_R) = \alpha_s(0)\left[1 - \alpha_s(0)\beta_0 t_R + (\alpha_s(0)\beta_0)^2 t_R^2 \right] + \order{t_R^3}
\end{equation}

Step 2: expand $\alpha_s^\text{NLO}(t_R)$.
Note that for the hint we use
\begin{equation}
    \ln(\alpha_s(0)/\alpha_s^\text{LO}(t_R)) = \ln(1 + \alpha_s(0)\beta_0 t_R) = \alpha_s(0) \beta_0 t_R - \frac 1 2 (\alpha_s(0)\beta_0)^2 t_R^2  + \order{t_R^3}
\end{equation}
by using \cref{eq:LOalphast}.
Together with step 1 we find
\begin{align}
    \alpha_s^\text{NLO}(t_R) &= \alpha_s(0)\left[1 - \alpha_s(0)\beta_0 t_R + (\alpha_s(0)\beta_0)^2 t_R^2 \right] \nonumber\\
    &\hspace{20pt} - \frac{\beta_1}{\beta_0} \alpha_s^2(0) \left[1 - 2\alpha_s(0)\beta_0 t_R + 3(\alpha_s(0)\beta_0)^2 t_R^2 \right]\nonumber\\
    &\hspace{70pt} \cdot \left[\alpha_s(0) \beta_0 t_R - \frac 1 2 (\alpha_s(0)\beta_0)^2 t_R^2 \right] + \order{t_R^3}\\
    &= \alpha_s(0)\left[1 + \left(-\beta_0\alpha_s(0) - \beta_1\alpha_s^2(0) \right) t_R\right. \nonumber\\
    &\hspace{50pt}\left. +\left(\beta_0^2 \alpha_s^2(0) + \frac 5 2 \beta_1 \beta_0 \alpha_s^3(0) \right)t_R^2 \right]  + \order{t_R^3} \label{eq:solresumNLLexp}
\end{align}
where the term $~\sim \alpha_s^3(0) t_R^2$ receives contributions from two summands.
Note that this is the full list of terms which appear up to $t_R^2$.

Step 3: Taylor-expand $\alpha_s^\text{NLO}(t_R)$.
First, we compute the second derivative:
\begin{align}
    \dv[2]{\alpha_s^\text{NLO}(t_R)}{t_R} &= \dv {t_R} \left( -\beta_0 \left(\alpha_s^\text{NLO}(t_R)\right)^2 -\beta_1 \left(\alpha_s^\text{NLO}(t_R)\right)^3 \right)\\
    &= -2\beta_0 \left(\alpha_s^\text{NLO}(t_R)\right) \dv{\alpha_s^\text{NLO}(t_R)}{t_R} -3\beta_1 \left(\alpha_s^\text{NLO}(t_R)\right)^2 \dv{\alpha_s^\text{NLO}(t_R)}{t_R}\\
    &= \left(-2\beta_0 \alpha_s^\text{NLO}(t_R) - 3\beta_1 \left(\alpha_s^\text{NLO}(t_R)\right)^2\right)\nonumber\\
    &\hspace{50pt} \cdot\left( -\beta_0 \left(\alpha_s^\text{NLO}(t_R)\right)^2 -\beta_1 \left(\alpha_s^\text{NLO}(t_R)\right)^3\right)\\
    &= 2\beta_0^2 \left(\alpha_s^\text{NLO}(t_R)\right)^3 + 5\beta_0\beta_1 \left(\alpha_s^\text{NLO}(t_R)\right)^4 + 3\beta_1^2 \left(\alpha_s^\text{NLO}(t_R)\right)^5
\end{align}
Now, we perform the Taylor-expansion:
\begin{align}
    \alpha_s^\text{NLO}(t_R) &= \alpha_s(0) + \frac 1 {1!}\left.\dv{\alpha_s^\text{NLO}(t_R)}{t_R}\right|_{t_R=0} t_R + \frac 1 {2!}\left.\dv[2]{\alpha_s^\text{NLO}(t_R)}{t_R}\right|_{t_R=0} t_R^2 + \order{t_R^3}\\
    &= \alpha_s(0)\left[1 + \left(-\beta_0\alpha_s(0) - \beta_1\alpha_s^2(0) \right) t_R\right. \nonumber\\
    &\hspace{50pt}\left. +\left(\beta_0^2 \alpha_s^2(0) + \frac 5 2 \beta_1 \beta_0 \alpha_s^3(0) + \frac 3 2 \beta_1^2 \left(\alpha_s^\text{NLO}(t_R)\right)^5 \right)t_R^2 \right]  + \order{t_R^3} \label{eq:solresumNLLTaylor}
\end{align}

Step 4: \cref{eq:solresumNLLexp} and \cref{eq:solresumNLLTaylor} agree at LL, $\alpha_s^k(0) t_R^{k}$ for $k=0,1,2$, and all terms are non-zero.

Step 5: \cref{eq:solresumNLLexp} and \cref{eq:solresumNLLTaylor} agree at NLL, $\alpha_s^k(0) t_R^{k-1}$ for $k=1,2,3$.
Note that the term $\alpha_s^1(0) t_R^{0}$ is trivially 0 in both expressions and thus typically not considered a part of the NLL resummation.

Step 6: \cref{eq:solresumNLLexp} and \cref{eq:solresumNLLTaylor} do not agree at NNLL, $\alpha_s^k(0) t_R^{k-2}$ for $k=2,3,4$.
Note that the terms $\alpha_s^2(0) t_R^{0}$ and $\alpha_s^3(0) t_R^{1}$ are both trivially 0 in both expressions and thus typically not considered a part of the NNLL resummation.
The first non-trivial term $\alpha_s^4(0) t_R^{2}$ appears in the Taylor-expansion \cref{eq:solresumNLLTaylor}, but not in the expansion of \cref{eq:solresumNLLexp}.
This implies that \cref{eq:NLOalphast} is not the true solution to \cref{eq:NLObetat}, as we say above.
However, \cref{eq:NLOalphast} does solve the beta function at the requested accuracy, \cref{box:NLOalphas}, and it does perform the requested NLL resummation.

Note that we merely prove here that the first few terms are available and not the full NLL accuracy, which requires that all terms $\alpha_s^k(0) t_R^{k-1}$ for $k=1,\cdots,\infty$ are present.
%The prove can be done using \cref{eq:LOalphasresum} and the series representation of the logarithm.
